\documentclass[11pt]{article}
\usepackage[a4paper,margin=1.1in]{geometry}
\usepackage{parskip}
\usepackage[hidelinks]{hyperref}
\pagestyle{empty}

\begin{document}

{\Large\bfseries Cover Letter}

\vspace{1em}

Dear Editor,

I would like to submit a manuscript titled ``PACE: Probabilistic Adaptive
Contention Control for Non-Primary Channel Access in IEEE 802.11bn'' by
Taewon Song and Yeunwoong Kyung for consideration in the IEEE Transactions
on Wireless Communications.

In this manuscript, we propose PACE, a distributed probabilistic contention
control scheme for non-primary channel access (NPCA) that lets stations
track the time-varying fair-share transmission probability of a depleting
transmission window without central coordination. The core problem is that
the IEEE 802.11bn draft has stations reuse binary exponential backoff on
the non-primary channel, yet that rule was designed for open-ended
contention: within a bounded window, the doubled backoff frequently
exceeds the time remaining, so stations fall silent while the channel
idles. PACE addresses this by combining a deadline-scaled initialization,
a multiplicative feedback rule in which stations copy the successful
sender's probability on solo receptions, and a self-exclusion rule that
silences stations whose frame no longer fits the remaining window. A
finite-horizon stochastic approximation analysis of the feedback rule
yields an inverse-square-root scaling law for the adaptation step in the
window length, and the resulting window-scaled variant, PACE-dynamic,
derives every design parameter of the scheme from the one quantity the
NPCA procedure already signals. Under IEEE 802.11 standard timing on a
channel shared with native stations, PACE keeps the total useful airtime
within 3\% of a fair-share reference under RTS/CTS-protected access,
exceeds the standard-compliant NPCA baseline by up to 61\% there and by up
to 4.1 times under basic access, and holds the visiting stations' airtime
at or below their population-proportional share, with PACE-dynamic
preserving this proportionality across visit durations where a fixed
coefficient drifts past it. An ablation study confirms that the
constituent mechanisms are necessary for these gains.

We believe this work fits the scope of the IEEE Transactions on Wireless
Communications in the design and analysis of wireless access protocols:
the analytical contribution characterizes distributed contention control
as a finite-horizon stochastic approximation over the bounded NPCA window,
and the protocol contribution is standard-compatible channel access for
next-generation Wi-Fi in dense overlapping-BSS deployments.

Both authors have reviewed and approved the manuscript. This manuscript
has not been published in any journal and is not under consideration for
publication elsewhere.

Sincerely yours,

Yeunwoong Kyung\\
Department of Electronic Engineering, Seoul National University of Science
and Technology\\
\href{mailto:ywkyung@seoultech.ac.kr}{ywkyung@seoultech.ac.kr}

\end{document}
